\documentclass[10pt]{article}
\usepackage[a4paper,margin=2.5cm,top=3cm,bottom=3cm]{geometry}
\usepackage{amsmath,amssymb}
\usepackage{booktabs}
\usepackage{enumitem}
\usepackage{xcolor}
\usepackage{graphicx}
\usepackage{tikz}
\usetikzlibrary{arrows.meta,positioning,shapes.geometric,calc}
\usepackage[colorlinks=true,linkcolor=blue,urlcolor=blue,citecolor=blue]{hyperref}
\usepackage{caption}
\usepackage{subcaption}
% 中文支持（需要 XeLaTeX 编译）
\usepackage{fontspec}
\setmainfont[Path=/System/Library/Fonts/,Extension=.ttc,Index=2]{PingFang.ttc}
\setsansfont[Path=/System/Library/Fonts/,Extension=.ttc,Index=2]{PingFang.ttc}

\newcommand{\techreportnumber}{TR-2026-09-ZH}

\begin{document}

% ── 封面页（斯坦福技术报告风格）──────────────────────
\begin{titlepage}
  \centering
  \vspace*{1.2cm}

  {\large\textsc{技术报告}}\\[0.3em]
  \rule{\linewidth}{1pt}\\[0.5em]

  {\Huge\bfseries
    Apple Silicon 上本地 LLM 推理服务器的自愈看门狗}\\[0.6em]

  {\Large oMLX 故障模式与有界恢复的实证研究}\\[1.2em]

  {\large ricky8848}\\[0.3em]
  {\normalsize 独立研究员}\\[2em]

  \begin{tabular}{rl}
    报告编号: & \techreportnumber \\[0.3em]
    日期:     & 2026年9月\\[0.3em]
    许可证:   & MIT \\[0.3em]
    代码仓库: & \url{https://github.com/ricky8848/omlx-watchdog} (v6.1)
  \end{tabular}\\[2em]

  {\small\itshape
    配套配置: \url{https://github.com/ricky8848/mac-m5-128g-omlx-settings}}

  \vfill
\end{titlepage}

% ── 摘要 ───────────────────────────────────────
\begin{abstract}
\noindent Apple Silicon 上的本地大语言模型（LLM）推理服务器越来越多地作为自主编码智能体的后端，但它们运行在任何编排器之外：没有 Kubernetes 存活探针、没有 systemd 重启策略——更关键的是，故障模式是\emph{部分性的}（进程存活但引擎卡死）。我们对一台这样的服务器进行了实证研究：\texttt{oMLX} 在 Apple M5 Max（128GB）上服务 Qwen3.8-27B，连续 13 天、共 18{,}516 次看门狗心跳。我们形式化了六种故障类别（引擎卡死、端口占用者、SIGKILL/OOM、陈旧注册表、预填充内存保护拒绝、客户端重试放大），并提出了一个无自干扰的健康模型：四个互补检查的谓词互不冲突，驱动有界升级阶梯（\texttt{omlx restart} $\rightarrow$ 仅监听者端口清理 + 30秒宽限 $\rightarrow$ GUI 重启）。在生产环境中，看门狗在一个 60 秒心跳内检测到所有引擎故障，并在四分钟内恢复服务且无需人工干预；部署后主导的故障类别（模型切换后的陈旧注册表）通过版本化修复被消除。我们讨论了后端自愈与客户端有界重试的正交性，以及投机解码吞吐量（MTP 接受率 85--96\%）与可用性的独立但乘积关系。

\vspace{1em}
\noindent\textbf{关键词：} 本地 LLM 推理、自愈系统、Apple Silicon、看门狗、故障分类学、有界升级
\end{abstract}

% ── 目录 ───────────────────────────────────────
\tableofcontents
\newpage

% ═══════════════════════════════════════════════════════════════════
\section{引言}\label{sec:intro}

自主编码智能体（DeepSeek Harness、Codex CLI、OpenClaw/Hermes 网关）向本地 LLM 服务器发起长时程工具调用工作负载。在 Apple Silicon 上，\texttt{oMLX} 通过 Metal/ANE 提供 OpenAI 兼容 API，支持多 token 预测（MTP）。与云部署不同，本地服务器\emph{没有编排器}：卡死的引擎（进程存活、token 冻结）对 \texttt{ps} 不可见，而朴素的"日志 mtime = 忙碌"存活探针在取消请求写入自身日志行时会自干扰。

本文贡献：
\begin{itemize}[nosep]
  \item 在 13 天生产跟踪中观察到的六种故障类别分类学（第~\ref{sec:taxonomy} 节）。
  \item 四个检查谓词互不冲突的无自干扰健康模型（第~\ref{sec:detection} 节）。
  \item 带形式化端口清理安全论证的有界升级阶梯（第~\ref{sec:recovery} 节）。
  \item 生产评估：18{,}516 次心跳下 MTTD $\le$ 60s、MTTR $<$ 4min（第~\ref{sec:eval} 节）。
  \item 结合服务器侧看门狗、客户端有界重试（$N{=}40$）、launchd 双重守护的三层防御体系（第~\ref{sec:architecture} 节）。
\end{itemize}

% ═══════════════════════════════════════════════════════════════════
\section{系统架构}\label{sec:architecture}

图~\ref{fig:arch} 展示了三层防御体系。第 1 层（launchd）提供粗粒度 \texttt{StartInterval=60} 心跳。第 2 层（看门狗脚本）评估四个谓词并通过有界阶梯升级。第 3 层（客户端重试预算）确保在途请求跨越重启窗口而不放大负载。

\begin{figure}[t]\centering
\begin{tikzpicture}[
  box/.style={draw, rounded corners=3pt, minimum width=4.2cm, minimum height=1.6cm,
              align=center, font=\small},
  arrow/.style={-{Stealth[length=3mm]}, thick}
]
\node[box, fill=blue!8] (launchd) {第 1 层: launchd\\ \texttt{StartInterval=60}};
\node[box, fill=green!8, below=1.2cm of launchd] (watchdog) {第 2 层: 看门狗\\ 4-检查谓词模型};
\node[box, fill=orange!8, below=1.2cm of watchdog] (ladder) {有界阶梯\\ restart $\to$ port-clean $\to$ GUI};
\node[box, fill=red!8, below=1.2cm of ladder] (client) {第 3 层: 客户端重试\\ \texttt{maxRetries=40}};

\draw[arrow] (launchd) -- node[right, font=\small]{心跳} (watchdog);
\draw[arrow] (watchdog) -- node[right, font=\small]{升级} (ladder);
\draw[arrow] (watchdog) -- node[left, font=\small]{健康} ++(-2.5,-0.8);

\draw[<->, arrow] (client) -- node[right, font=\small]{跨越重启窗口} ++(2.5,-0.8);
\end{tikzpicture}
\caption{三层防御：launchd 心跳 $\to$ 看门狗谓词评估 $\to$ 有界升级阶梯，客户端重试预算跨越重启窗口。}
\label{fig:arch}
\end{figure}

% ═══════════════════════════════════════════════════════════════════
\section{故障分类学}\label{sec:taxonomy}

表~\ref{tab:fails} 列出了六种故障类别及其生产发生率。图~\ref{fig:taxonomy} 展示了故障类别与恢复动作之间的因果关系。

\begin{table}[t]\centering
\small
\renewcommand{\arraystretch}{1.3}
\begin{tabular}{@{}lll@{}}
\toprule
类别 & 机制 & 发生率（13天） \\
\midrule
F1 引擎卡死   & Metal 层卡住；请求被接受但无 token & 2 次事件（9月8日） \\
F2 端口占用者  & 重启后 :8000 上陈旧监听；重试旋转 & 26 次端口清理 \\
F3 SIGKILL/OOM    & 系统杀死；launchd 不自动重启 GUI 应用 & 96 次回退 \\
F4 陈旧注册表 & \texttt{loaded\_models} 在模型切换后不同步 & 79 次恢复（v5 时代） \\
F5 预填充保护拒绝 & 内存上限处的可预测 HTTP 400（设计使然） & n/a \\
F6 客户端重试放大 & \texttt{mode:always} 重试 $\times$ 重启窗口 = 死亡循环（修复前） & 1 次事件 \\
\bottomrule
\end{tabular}
\caption{故障类别与生产发生率（9月27日 -- 9月9日，18{,}516 次心跳）。F4 主导 v5 时代，通过 v6 的空注册表感知检查 B 被消除。}
\label{tab:fails}
\end{table}

\begin{figure}[t]\centering
\begin{tikzpicture}[
  fbox/.style={draw, rounded corners=2pt, minimum width=3.5cm, minimum height=1.4cm,
               align=center, font=\small},
  rbox/.style={draw, rounded corners=2pt, minimum width=3.5cm, minimum height=1.4cm,
               align=center, font=\small, fill=green!5},
  arrow/.style={-{Stealth[length=3mm]}, thick}
]
\node[fbox, fill=red!8] (F1) {F1: 引擎卡死};
\node[fbox, fill=red!8, right=1.5cm of F1] (F2) {F2: 端口占用者};
\node[fbox, fill=red!8, below=1.2cm of F1] (F3) {F3: SIGKILL/OOM};
\node[fbox, fill=red!8, right=1.5cm of F3] (F4) {F4: 陈旧注册表};

\node[rbox, below=1.5cm of F3] (R1) {阶梯步骤 1:\\ \texttt{omlx restart}};
\node[rbox, right=1.5cm of R1] (R2) {阶梯步骤 2:\\ PortClean + kill};
\node[rbox, below=1.5cm of R2] (R3) {阶梯步骤 3:\\ \texttt{open -a oMLX}};

\draw[arrow] (F1) -- node[above, font=\tiny]{卡死} (R1);
\draw[arrow] (F2) -- node[right, font=\tiny]{陈旧监听} (R2);
\draw[arrow] (F3) -- node[left, font=\tiny]{进程死亡} (R1);
\draw[arrow] (F4) -- node[right, font=\tiny]{注册表为空} (R1);
\draw[arrow] (R2) -- node[right, font=\tiny]{若重启失败} (F1);
\draw[arrow] (R3) -- node[left, font=\tiny]{最后手段} ++(-1.5,-0.8);
\end{tikzpicture}
\caption{故障类别到恢复阶梯步骤的因果映射。F2（端口占用者）是唯一需要步骤 2 的类别；其余均由步骤 1 解决或通过版本化修复消除。}
\label{fig:taxonomy}
\end{figure}

F1/F3 的区分很重要：\texttt{oMLX restart} 修复卡死，但被占用的端口使重启本身失败，产生 F2 旋转——这正是仅监听者清理的动机。

% ═══════════════════════════════════════════════════════════════════
\section{检测：无自干扰健康模型}\label{sec:detection}

看门狗心跳 $t$ 评估服务器状态的四个谓词：
\begin{align}
A(t) &: \texttt{/api/status 在 5s 内可达} \\[0.3em]
B(t) &: \big(\texttt{loaded\_models} = \varset\big) \lor \big(\exists m{:}\; m.\texttt{registered}\big) \\[0.3em]
C(t) &: \neg\,\textsf{busy}(t) \;\lor\; \textsc{IdleProbe}(\texttt{max\_tokens}=4) = \textsc{ok} \\[0.3em]
D(t) &: \neg\,\big(\textsf{workInFlight}(t) \land \Delta t_{\text{token}}(t) \ge 900\,\text{s}\big)
\end{align}

其中 $\textsf{busy}(t)$ 是服务器自身的队列计数器（\texttt{active+waiting}$>$0）。健康谓词为 $H(t)=A\land B\land C\land D$。

\noindent\textbf{无自干扰性。} 检查 C \emph{仅在队列可证明为空时}触发；长时间解码（数小时工具调用）保持 \texttt{active}$\ge 1$，永远不会被探测。检查 D 要求工作\emph{在途}且 token 计数器冻结——空闲服务器没有可冻结的引用。v2 设计（日志 mtime）违反了这一点：取消请求写入日志行，将死引擎标记为"忙碌"。

\noindent\textbf{空注册表感知（F4 修复）。} 模型切换后注册表可能暂时为空；$B(t)$ 将 $\varset$ 视为健康而非触发重启风暴（v5 时代有 79 次此类事件）。

图~\ref{fig:detection} 展示了单次心跳的决策流程。

\begin{figure}[t]\centering
\begin{tikzpicture}[
  dbox/.style={draw, rounded corners=2pt, minimum width=4cm, minimum height=1.5cm,
               align=center, font=\small},
  ddec/.style={draw, diamond, aspect=2.5, minimum width=3cm, align=center, font=\small},
  arrow/.style={-{Stealth[length=3mm]}, thick}
]
\node[dbox, fill=blue!5] (start) {心跳 $t$: 读取\\ \texttt{/api/status}};
\node[ddec, below=1.2cm of start] (A) {$A(t)$?};
\node[dbox, below right=1.2cm and 0.5cm of A] (B) {检查 $B(t)$:\\ 注册表状态};
\node[dbox, below=1.2cm of B] (C) {检查 $C(t)$:\\ 空闲探测};
\node[dbox, below=1.2cm of C] (D) {检查 $D(t)$:\\ token 停滞};
\node[dbox, below=1.2cm of D] (H) {$H(t)=A\land B\land C\land D$};
\node[dbox, below right=1.2cm and 0.5cm of H, fill=green!8] (ok) {健康\\ 记录心跳};
\node[dbox, below left=1.2cm and 0.5cm of H, fill=red!8] (recover) {恢复\\ 升级阶梯};

\draw[arrow] (start) -- (A);
\draw[arrow, dashed] (A) -- node[right, font=\tiny]{不可达} ++(1.5,-0.3) -- (recover);
\draw[arrow] (A) -- node[left, font=\tiny]{可达} (B);
\draw[arrow] (B) -- node[right, font=\tiny]{空或已注册} (C);
\draw[arrow] (B) -- node[right, font=\tiny]{不同步} ++(1.5,-0.3) -- (recover);
\draw[arrow] (C) -- node[right, font=\tiny]{空闲或探测 ok} (D);
\draw[arrow] (C) -- node[right, font=\tiny]{探测失败} ++(1.5,-0.3) -- (recover);
\draw[arrow] (D) -- node[right, font=\tiny]{无停滞} (H);
\draw[arrow] (D) -- node[right, font=\tiny]{停滞 $\ge$ 900s} ++(1.5,-0.3) -- (recover);
\draw[arrow] (H) -- node[right, font=\tiny]{真} (ok);
\draw[arrow] (H) -- node[left, font=\tiny]{假} ++(-1.5,-0.3) -- (recover);
\end{tikzpicture}
\caption{单次看门狗心跳的决策流程。每个检查仅在前序全部通过时评估；任何失败路由到有界升级阶梯（第~\ref{sec:recovery} 节）。}
\label{fig:detection}
\end{figure}

% ═══════════════════════════════════════════════════════════════════
\section{恢复：有界升级阶梯}\label{sec:recovery}

当 $H(t)=\textsc{false}$ 时，看门狗升级：
\begin{enumerate}[nosep]
  \item $\texttt{omlx restart --timeout 240}$（引擎级，不杀进程）
  \item 失败时：\textsc{PortClean} -- 用 $\texttt{lsof -sTCP:LISTEN}$ 在 :8000 上识别 PID，\texttt{kill}，30秒宽限，然后 \texttt{kill -9}；重试步骤 1
  \item 最后手段：$\texttt{open -a oMLX}$（GUI 重启）
\end{enumerate}

PID 锁文件防止重入心跳。\textbf{PortClean 安全性：} kill 目标被验证为端口 8000 上的唯一 TCP 监听者——从不匹配进程列表——因此不可能向无关进程发送信号。阶梯是有界的：每次心跳最多一次重启、一次端口清理、一次 GUI 启动；下次心跳从干净状态重新评估。

图~\ref{fig:ladder} 展示了带时间边界的阶梯。

\begin{figure}[t]\centering
\begin{tikzpicture}[
  lbox/.style={draw, rounded corners=2pt, minimum width=5cm, minimum height=1.6cm,
               align=center, font=\small},
  arrow/.style={-{Stealth[length=3mm]}, thick}
]
\node[lbox, fill=blue!5] (step1) {步骤 1: \texttt{omlx restart --timeout 240}\\[0.3em]
  {\small 引擎级重置，不杀进程}};
\node[lbox, fill=orange!8, below=1.5cm of step1] (step2) {步骤 2: \textsc{PortClean}\\[0.3em]
  {\small lsof :8000 $\to$ kill PID, 30s 宽限, 重试步骤 1}};
\node[lbox, fill=red!8, below=1.5cm of step2] (step3) {步骤 3: \texttt{open -a oMLX}\\[0.3em]
  {\small GUI 重启（最后手段）}};

\draw[arrow] (step1) -- node[right, font=\small]{超时 240s / 失败} (step2);
\draw[arrow] (step2) -- node[right, font=\small]{端口仍被占用} (step3);

\node[font=\tiny, right=0.5cm of step1] {MTTR 贡献: $\sim$30s};
\node[font=\tiny, right=0.5cm of step2] {MTTR 贡献: $\sim$60s};
\node[font=\tiny, right=0.5cm of step3] {MTTR 贡献: $\sim$90s};
\end{tikzpicture}
\caption{带时间边界的有界升级阶梯。每步每次心跳最多尝试一次；下次心跳从干净状态重新评估。}
\label{fig:ladder}
\end{figure}

% ═══════════════════════════════════════════════════════════════════
\section{评估：13 天生产数据}\label{sec:eval}

\noindent\textbf{设置。} Apple M5 Max，128GB 统一内存；macOS Tahoe；oMLX v0.6.x 服务 Qwen3.8-27B（oQ8e，后 oQ4e）+ MTP；launchd \texttt{StartInterval=60}。客户端：DSH Web GUI、OpenClaw/Hermes 网关（同一端点）。

\begin{table}[t]\centering
\small
\renewcommand{\arraystretch}{1.3}
\begin{tabular}{@{}lrr@{}}
\toprule
指标 & 数值 \\
\midrule
总心跳数 / 挂钟时间          & 18{,}516 / 13.2 天 \\
HEALTHY 空闲心跳               & 13{,}527 (73.0\%) \\
BUSY 心跳                       & 4{,}188 (22.6\%) \\
恢复事件完成        & 79 RECOVERED, 0 STILL\_DOWN（阶梯后） \\
平均检测时间 (MTTD)       & $\le 60$\,s（一次心跳） \\
平均恢复时间 (MTTR)      & $<$ 4\,min（重启 + 重载观察） \\
误报重启 (v6 时代) & 最后 48h 窗口内 0 \\
解码吞吐量 (BUSY 心跳)   & 中位 1.9 tok/s; p90 7.6 tok/s（长上下文, TQKV+MTP） \\
\bottomrule
\end{tabular}
\caption{生产统计，\texttt{watchdog.log}, 2026-08-27 $\rightarrow$ 2026-09-09。}
\label{tab:eval}
\end{table}

9 月 8 日窗口（F1 卡死 + F4 陈旧注册表，v6 前）展示了阶梯在高负载下的表现：96 次回退升级全部完成恢复，v4 端口清理修复消除了 F2 旋转。v6 之后（9 月 9 日），48h 内零误报重启。

\noindent\textbf{吞吐量上下文。} oQ4e + MTP（接受率 85--96\%）+ TurboQuant Q4 KV 缓存下，短上下文解码达 $\sim$30--50 tok/s；长上下文（80k+ token）稳定在 12--25 tok/s，与 M5 Max $\sim$614\,GB/s 内存带宽对权重+KV 读取的约束一致。可用性（本文）与吞吐量正交：三层防御体系将二者相乘。

图~\ref{fig:timeline} 展示了 13 天跟踪的事件时间线。

\begin{figure}[t]\centering
\begin{tikzpicture}[xscale=0.35, yscale=1]
\draw[->] (0,-0.5) -- (26,0) node[right]{天};
\draw[->] (-0.5,-12) -- (0,1) node[left]{事件};

% Day 0-5: healthy
\draw[blue, thick] (1,-2) -- (4.5,-2);
\node[font=\tiny, above] at (2.75,-2) {HEALTHY};

% Day 6-8: F1 + F4 window
\draw[red, thick] (5,-2) -- (8.5,-2);
\node[font=\tiny, above] at (6.75,-2) {F1 卡死 + F4};
\node[font=\tiny, below] at (6.75,-2) {96 次回退};

% Day 8: v4 port-clean fix
\draw[orange, thick] (8.5,-2) -- (9,-2);
\node[font=\tiny, above] at (8.75,-2) {v4 修复};

% Day 9-13: v6 era, healthy
\draw[green!60!black, thick] (9,-2) -- (13.5,-2);
\node[font=\tiny, above] at (11.25,-2) {v6: 0 误报};

% Tick marks
\foreach \x/\lbl in {1/Day 0, 3.5/Day 4, 7/Day 8, 10.5/Day 12}
  \draw (\x,-0.3) -- (\x,0) node[below, font=\tiny]{\lbl};

% Event counts
\node[font=\small] at (6.75,0.8) {13{,}527 HEALTHY};
\node[font=\small] at (11.25,0.8) {4{,}188 BUSY};
\end{tikzpicture}
\caption{13 天跟踪的事件时间线。9 月 8 日窗口（F1+F4）展示了阶梯在高负载下的表现；v6 之后零误报重启。}
\label{fig:timeline}
\end{figure}

% ═══════════════════════════════════════════════════════════════════
\section{相关工作}\label{sec:related}

Kubernetes 存活/就绪探针假设有进程杀死权限的编排器；我们的阶梯是单主机类比，其中"编排器"（launchd）故意\emph{不}自动重启 GUI 应用。SRE 自愈模式（有界升级、无重试放大）直接映射到阶梯步骤 1--3 和客户端侧 \texttt{maxRetries=40} 边界。本地 LLM 工具（Ollama、llama.cpp server）未提供等效看门狗；DFlash/MTP 投机解码改善每请求吞吐量但不解决部分可用性问题。

% ═══════════════════════════════════════════════════════════════════
\section{效度威胁}

单一操作者、单台机器、一个服务器版本族。\texttt{loaded\_models} 字段是稳定 API 表面但非契约性的；检查 B 在重命名时优雅降级（空注册表感知）。吞吐量数字是跟踪推导的中位数，非受控基准测试。

% ═══════════════════════════════════════════════════════════════════
\section{结论}

一个 180 行 shell 脚本，四个非冲突谓词和一个有界阶梯，为无编排器的本地 LLM 服务器提供完整自愈：MTTD $\le 60$s、MTTR $<$4min，13 天智能体工作负载零人工干预。该设计通过替换重启原语推广到任何 OpenAI 兼容本地后端（llama.cpp、Ollama）。

% ═══════════════════════════════════════════════════════════════════
\section*{工件可用性}

看门狗 + 安装器：\url{https://github.com/ricky8848/omlx-watchdog}（MIT, v6.1; \texttt{curl | bash}）。三层配置：\url{https://github.com/ricky8848/mac-m5-128g-omlx-settings}。完整去标识化跟踪：\texttt{docs/paper/production-data.json}（在工件仓库中）。

\end{document}
